\section{\gls{etram} Dataset}
\label{sec:background_etram}

\subsection{Dataset Overview \& Acquisition Setup}

\gls{etram} is a first-of-its-kind, fully event-based traffic monitoring dataset introduced by \textcite{Verma_2024_CVPR} at CVPR 2024.
Data was captured using the Prophesee EVK4 HD event camera at $1280 \times 720$ spatial resolution with a temporal resolution exceeding $10{,}000$\,fps, mounted statically at approximately $6$\,m height with a $35^\circ$ pitch angle.
Unlike ego-motion datasets, where the camera itself moves through the scene, this dataset provides a static perspective; a fixed vantage point ensures that events are generated predominantly by moving traffic participants rather than background infrastructure.
Recordings span intersections, roadways, and local streets around Arizona State University, Tempe campus, totalling ten hours of annotated event data.


\subsection{Class Taxonomy, Scale \& Annotation}

\gls{etram} covers eight distinct object classes organised into three broad categories: Vehicles (Car, Truck, Bus, Tram), Pedestrians, and Micro-mobility (Bicycle, Bike, Wheelchair).
Over $2$M bounding boxes were manually annotated using CVAT \cite{Verma_2024_CVPR}, with each annotation additionally carrying an object ID to support multi-object tracking.
Raw data is distributed in multiple formats (RAW, DAT, and H5), with ground-truth annotations provided in NumPy format.
For this project, the dataset was abstracted to the Vehicle category exclusively, allowing the \gls{snn}--\gls{cnn} comparison to be isolated on a single, well-defined object class.

\subsection{Suitability for This Project}

Two properties motivated the selection of \gls{etram} for this project.
First, critically for an \gls{adas} and autonomous driving context, recordings span the full spectrum of real-world operating conditions: daytime, nighttime, and twilight, across sunny, overcast, and rainy weather - precisely the range of environments an autonomous perception system must handle reliably, and one that deliberately exercises the conditions under which frame-based sensors are known to degrade.
Second, the \gls{etram} study provides a dataset-specific spatiotemporal filtering and denoising pipeline tuned to the noise characteristics of the EVK4 sensor, which was adopted directly as the signal conditioning stage of the \texttt{REPLICA} preprocessing system, This eliminated the need to develop bespoke noise suppression from first principles, and significantly accelerated the transition to model development.
The full \texttt{REPLICA} pipeline is detailed in Section~\ref{sec:data_preprocessing}.
